\input format.tex

\usepackage{graphicx}
\graphicspath{{cores/}}

\usepackage{environ}
\usepackage{colortbl,array,booktabs}
\usepackage{tabularx}

\colorlet{TablaBordeSuperior}{topcolor}
\colorlet{TablaBordeInferior}{topcolor}
\colorlet{TablaCentroSuperior}{blue!1}
\colorlet{TablaCentroInferior}{blue!20}
\colorlet{FuenteCabeceraTabla}{white}

\newcolumntype{M}[1]{>{\centering\arraybackslash}m{#1}}

\tcbset{rtab/.style={
freelance,
frame code={
 \path[top color=topcolor,bottom color=topcolor]
   ([yshift=-#1*(\baselineskip+2pt)]interior.north west) --
   ([yshift=-#1*(\baselineskip+2pt)]interior.north east) {[sharp corners]--
    ([yshift=3pt]interior.north east) --
    ([yshift=3pt]interior.north west)} -- cycle;

  },
interior code={},
 }
}

\newcommand\fuentecabecera[1]{\textcolor{black}{\textbf{#1}}}

\NewEnviron{RCtable}[1]
  {%
    \begin{tcolorbox}[#1]\BODY\end{tcolorbox}%
  }



\begin{document}

\vspace*{2mm}
%% 各章节
\setlength{\arrayrulewidth}{1pt}
\fontsize{9.3pt}{11pt}\selectfont
\color{gray2}

{\noindent\bf\sanhao 二、肠道菌含量}

\vspace*{2mm}

\begin{LRaside}[.8]{检测说明}
\noindent
\includegraphics[width=\linewidth]{hanliang.pdf}
\asidebreak %
正常人体肠道内寄居着数量庞大、种类繁多的微生物，其中以细菌为主，包括有益菌和有害菌。一般情况下，这些细菌的种类和数量都是相对固定的，它们在肠道内各司其职，共同维持肠道的正常运转。另外，食物中含有不常驻肠道的外源性致病菌，若不慎摄入，则有可能在肠道内大量繁殖，导致相关疾病。
\end{LRaside}

\vspace*{3mm}
\noindent
2.1 您的常见肠道菌含量
\vspace*{1.5mm}

%\begin{RCtable}{
%  rtab=1,
%  tabularx*={\arrayrulecolor{gray2}}%
%    {M{0.3\linewidth} M{0.31\linewidth}M{0.34\linewidth}}
%}
%\fuentecabecera{测试项目} & \fuentecabecera{健康状态提示} & \fuentecabecera{检测结果评价} \\[3pt]
%\bf{测试项目} & \bf{健康状态提示} & \bf{检测结果评价} \\[3pt]
%\end{RCtable}

\begin{tctabularx}{tabularx={m{4.8cm}<{\centering}m{5.2cm}<{\centering}m{4.55cm}<{\centering}}}
&&
\\[-6pt]
\cellcolor{topcolor} \raisebox{2.614pt}{\color{black00}\bf 检测项目} &
\cellcolor{topcolor} \raisebox{2.614pt}{\color{black00}\bf 健康状态提示} &
\cellcolor{topcolor} \raisebox{2.614pt}{\color{black00}\bf 检测结果评价}
\end{tctabularx}

\vspace*{-4.25mm}
\fontsize{8.8pt}{11pt}\selectfont
\arrayrulecolor{gray2}
\begin{longtable}{m{4.8cm}m{5.2cm}<{\centering}m{0cm}@{}m{4.61cm}<{\centering}}
\hline
\parbox[c]{\hsize}{\vskip7pt {\lantxh 拟杆菌属\\常见肠道菌，多数与人体互利共生，少数菌种可能会导致菌血症等} \vskip7pt} & \parbox[c]{\hsize}{\vskip7pt\centerline{\raisebox{-1.5ex}{\includegraphics[width=5cm,height=0.55cm]{bar.pdf}}}\vskip7pt}  &
\hspace*{-1.51cm}\raisebox{-0.45ex}{\includegraphics{cry.pdf}}
 & \begin{minipage}{4.60cm}\begin{center}{{\color{red}\lantxh 高{\\ \bahao 不利于肠道菌群平衡}} }\end{center} \end{minipage} \\
% & \begin{minipage}{4.60cm}\begin{center}{{\color{red}\lantxh 高\\ \bahao 不利于肠道菌群平衡} }\end{center} \end{minipage} \\
\hline
\parbox[c]{\hsize}{\vskip7pt {\lantxh 布劳特氏菌属\\发酵多种植物多糖产生乙酸盐，促进肠道健康} \vskip7pt} & \parbox[c]{\hsize}{\vskip7pt\centerline{\raisebox{-1.5ex}{\includegraphics[width=5cm,height=0.55cm]{bar.pdf}}}\vskip7pt}  &
\hspace*{-4.83cm}\raisebox{-0.45ex}{\includegraphics{smile.pdf}}
 & \begin{minipage}{4.60cm}\begin{center}{{\lantxh 高{\\ \bahao 有利于产生有益物质及肠道健康}} }\end{center} \end{minipage} \\
% & \begin{minipage}{4.60cm}\begin{center}{{\color{green}\lantxh 高\\ \bahao 有利于产生有益物质及肠道健康} }\end{center} \end{minipage} \\
\hline
\parbox[c]{\hsize}{\vskip7pt {\lantxh 栖粪杆菌属\\发酵纤维素产生丁酸等有益物质，抑制肠道炎症，促进肠道健康} \vskip7pt} & \parbox[c]{\hsize}{\vskip7pt\centerline{\raisebox{-1.5ex}{\includegraphics[width=5cm,height=0.55cm]{bar.pdf}}}\vskip7pt}  &
\hspace*{-4.83cm}\raisebox{-0.45ex}{\includegraphics{smile.pdf}}
 & \begin{minipage}{4.60cm}\begin{center}{{\lantxh 高{\\ \bahao 有利于产生有益物质及抑制肠道炎症}} }\end{center} \end{minipage} \\
% & \begin{minipage}{4.60cm}\begin{center}{{\color{green}\lantxh 高\\ \bahao 有利于产生有益物质及抑制肠道炎症} }\end{center} \end{minipage} \\
\hline
\parbox[c]{\hsize}{\vskip7pt {\lantxh 粪球菌属\\发酵多糖，能够产生丁酸、乙酸等有益物质，调节肠道环境} \vskip7pt} & \parbox[c]{\hsize}{\vskip7pt\centerline{\raisebox{-1.5ex}{\includegraphics[width=5cm,height=0.55cm]{bar.pdf}}}\vskip7pt}  &
\hspace*{-4.83cm}\raisebox{-0.45ex}{\includegraphics{smile.pdf}}
 & \begin{minipage}{4.60cm}\begin{center}{{\lantxh 高{\\ \bahao 有利于产生有益物质及调节肠道环境}} }\end{center} \end{minipage} \\
% & \begin{minipage}{4.60cm}\begin{center}{{\color{green}\lantxh 高\\ \bahao 有利于产生有益物质及调节肠道环境} }\end{center} \end{minipage} \\
\hline
\parbox[c]{\hsize}{\vskip7pt {\lantxh 瘤胃球菌属\\帮助降解纤维素等人体不能消化的多糖，在肠道中广泛分布} \vskip7pt} & \parbox[c]{\hsize}{\vskip7pt\centerline{\raisebox{-1.5ex}{\includegraphics[width=5cm,height=0.55cm]{bar.pdf}}}\vskip7pt}  &
\hspace*{-4.83cm}\raisebox{-0.45ex}{\includegraphics{smile.pdf}}
 & \begin{minipage}{4.60cm}\begin{center}{{\lantxh 高{\\ \bahao 有利于降解纤维素等多糖}} }\end{center} \end{minipage} \\
% & \begin{minipage}{4.60cm}\begin{center}{{\color{green}\lantxh 高\\ \bahao 有利于降解纤维素等多糖} }\end{center} \end{minipage} \\
\hline
\parbox[c]{\hsize}{\vskip7pt {\lantxh 颤螺菌属\\帮助抗性淀粉和脂肪消化，保持正常体重，抑制肠道炎症} \vskip7pt} & \parbox[c]{\hsize}{\vskip7pt\centerline{\raisebox{-1.5ex}{\includegraphics[width=5cm,height=0.55cm]{bar.pdf}}}\vskip7pt}  &
\hspace*{-4.83cm}\raisebox{-0.45ex}{\includegraphics{smile.pdf}}
 & \begin{minipage}{4.60cm}\begin{center}{{\lantxh 高{\\ \bahao 有利于保持正常体重及抑制肠道炎症}} }\end{center} \end{minipage} \\
% & \begin{minipage}{4.60cm}\begin{center}{{\color{green}\lantxh 高\\ \bahao 有利于保持正常体重及抑制肠道炎症} }\end{center} \end{minipage} \\
\hline
\parbox[c]{\hsize}{\vskip7pt {\lantxh 副拟杆菌属\\帮助消化纤维素、抗性淀粉，保护肠道，抑制肠道炎症} \vskip7pt} & \parbox[c]{\hsize}{\vskip7pt\centerline{\raisebox{-1.5ex}{\includegraphics[width=5cm,height=0.55cm]{bar.pdf}}}\vskip7pt}  &
\hspace*{-4.83cm}\raisebox{-0.45ex}{\includegraphics{smile.pdf}}
 & \begin{minipage}{4.60cm}\begin{center}{{\lantxh 高{\\ \bahao 有利于抑制肠道炎症及保护肠道}} }\end{center} \end{minipage} \\
% & \begin{minipage}{4.60cm}\begin{center}{{\color{green}\lantxh 高\\ \bahao 有利于抑制肠道炎症及保护肠道} }\end{center} \end{minipage} \\
\hline
\parbox[c]{\hsize}{\vskip7pt {\lantxh 毛螺菌属\\发酵多种糖类产生乙酸、甲酸等物质，能保护肠黏膜，抑制肠道炎症} \vskip7pt} & \parbox[c]{\hsize}{\vskip7pt\centerline{\raisebox{-1.5ex}{\includegraphics[width=5cm,height=0.55cm]{bar.pdf}}}\vskip7pt}  &
\hspace*{-4.83cm}\raisebox{-0.45ex}{\includegraphics{smile.pdf}}
 & \begin{minipage}{4.60cm}\begin{center}{{\lantxh 高{\\ \bahao 有利于产生有益物质及抑制肠道炎症}} }\end{center} \end{minipage} \\
% & \begin{minipage}{4.60cm}\begin{center}{{\color{green}\lantxh 高\\ \bahao 有利于产生有益物质及抑制肠道炎症} }\end{center} \end{minipage} \\
\hline
\parbox[c]{\hsize}{\vskip7pt {\lantxh 多尔氏菌属\\肠道的主要产气菌之一，与肠易激综合征等疾病相关} \vskip7pt} & \parbox[c]{\hsize}{\vskip7pt\centerline{\raisebox{-1.5ex}{\includegraphics[width=5cm,height=0.55cm]{bar.pdf}}}\vskip7pt}  &
\hspace*{-1.51cm}\raisebox{-0.45ex}{\includegraphics{cry.pdf}}
 & \begin{minipage}{4.60cm}\begin{center}{{\color{red}\lantxh 高{\\ \bahao 可能增加肠易激综合征风险}} }\end{center} \end{minipage} \\
% & \begin{minipage}{4.60cm}\begin{center}{{\color{red}\lantxh 高\\ \bahao 可能增加肠易激综合征风险} }\end{center} \end{minipage} \\
\hline
\parbox[c]{\hsize}{\vskip7pt {\lantxh 链球菌属\\多数为共生菌，少数在免疫力弱时可能引起化脓性炎症、肺炎等} \vskip7pt} & \parbox[c]{\hsize}{\vskip7pt\centerline{\raisebox{-1.5ex}{\includegraphics[width=5cm,height=0.55cm]{bar.pdf}}}\vskip7pt}  &
\hspace*{-1.51cm}\raisebox{-0.45ex}{\includegraphics{cry.pdf}}
 & \begin{minipage}{4.60cm}\begin{center}{{\color{red}\lantxh 高{\\ \bahao 可能增加化脓性炎症、肺炎等的风险}} }\end{center} \end{minipage} \\
% & \begin{minipage}{4.60cm}\begin{center}{{\color{red}\lantxh 高\\ \bahao 可能增加化脓性炎症、肺炎等的风险} }\end{center} \end{minipage} \\
\hline
\parbox[c]{\hsize}{\vskip7pt {\lantxh 普雷沃氏菌属\\肠道共生菌，偏素食人群中含量高，少数菌种与感染、牙周疾病有关} \vskip7pt} & \parbox[c]{\hsize}{\vskip7pt\centerline{\raisebox{-1.5ex}{\includegraphics[width=5cm,height=0.55cm]{bar.pdf}}}\vskip7pt}  &
\hspace*{-1.51cm}\raisebox{-0.45ex}{\includegraphics{cry.pdf}}
 & \begin{minipage}{4.60cm}\begin{center}{{\color{red}\lantxh 高{\\ \bahao 不利于肠道菌群平衡}} }\end{center} \end{minipage} \\
% & \begin{minipage}{4.60cm}\begin{center}{{\color{red}\lantxh 高\\ \bahao 不利于肠道菌群平衡} }\end{center} \end{minipage} \\
\hline
\parbox[c]{\hsize}{\vskip7pt {\lantxh 罗斯拜瑞氏菌属\\产生丁酸等有益物质，抑制肠道炎症，有利于肠道及人体健康} \vskip7pt} & \parbox[c]{\hsize}{\vskip7pt\centerline{\raisebox{-1.5ex}{\includegraphics[width=5cm,height=0.55cm]{bar.pdf}}}\vskip7pt}  &
\hspace*{-4.83cm}\raisebox{-0.45ex}{\includegraphics{smile.pdf}}
 & \begin{minipage}{4.60cm}\begin{center}{{\lantxh 高{\\ \bahao 有利于产生有益物质及抑制肠道炎症}} }\end{center} \end{minipage} \\
% & \begin{minipage}{4.60cm}\begin{center}{{\color{green}\lantxh 高\\ \bahao 有利于产生有益物质及抑制肠道炎症} }\end{center} \end{minipage} \\
\hline
\parbox[c]{\hsize}{\vskip7pt {\lantxh 萨特氏菌属\\可能与胃肠道感染相关} \vskip7pt} & \parbox[c]{\hsize}{\vskip7pt\centerline{\raisebox{-1.5ex}{\includegraphics[width=5cm,height=0.55cm]{bar.pdf}}}\vskip7pt}  &
\hspace*{-1.51cm}\raisebox{-0.45ex}{\includegraphics{cry.pdf}}
 & \begin{minipage}{4.60cm}\begin{center}{{\color{red}\lantxh 高{\\ \bahao 可能增加胃肠道感染风险}} }\end{center} \end{minipage} \\
% & \begin{minipage}{4.60cm}\begin{center}{{\color{red}\lantxh 高\\ \bahao 可能增加胃肠道感染风险} }\end{center} \end{minipage} \\
\hline
\parbox[c]{\hsize}{\vskip7pt {\lantxh 双歧杆菌属\\有益菌，降解人体不能消化的多糖，产乳酸，调节免疫及肠道环境} \vskip7pt} & \parbox[c]{\hsize}{\vskip7pt\centerline{\raisebox{-1.5ex}{\includegraphics[width=5cm,height=0.55cm]{bar.pdf}}}\vskip7pt}  &
\hspace*{-4.83cm}\raisebox{-0.45ex}{\includegraphics{smile.pdf}}
 & \begin{minipage}{4.60cm}\begin{center}{{\lantxh 高{\\ \bahao 有利于调节免疫及肠道环境}} }\end{center} \end{minipage} \\
% & \begin{minipage}{4.60cm}\begin{center}{{\color{green}\lantxh 高\\ \bahao 有利于调节免疫及肠道环境} }\end{center} \end{minipage} \\
\hline
\parbox[c]{\hsize}{\vskip7pt {\lantxh 阿克曼氏菌属\\降解粘蛋白、调节免疫，有利于肠黏膜完整性，保持正常体重} \vskip7pt} & \parbox[c]{\hsize}{\vskip7pt\centerline{\raisebox{-1.5ex}{\includegraphics[width=5cm,height=0.55cm]{bar.pdf}}}\vskip7pt}  &
\hspace*{-4.83cm}\raisebox{-0.45ex}{\includegraphics{smile.pdf}}
 & \begin{minipage}{4.60cm}\begin{center}{{\lantxh 高{\\ \bahao 有利于调节免疫及保持正常体重}} }\end{center} \end{minipage} \\
% & \begin{minipage}{4.60cm}\begin{center}{{\color{green}\lantxh 高\\ \bahao 有利于调节免疫及保持正常体重} }\end{center} \end{minipage} \\
\hline
\parbox[c]{\hsize}{\vskip7pt {\lantxh 梭菌属\\多数为致病菌，可能引起腹泻、肠炎等疾病} \vskip7pt} & \parbox[c]{\hsize}{\vskip7pt\centerline{\raisebox{-1.5ex}{\includegraphics[width=5cm,height=0.55cm]{bar.pdf}}}\vskip7pt}  &
\hspace*{-1.51cm}\raisebox{-0.45ex}{\includegraphics{cry.pdf}}
 & \begin{minipage}{4.60cm}\begin{center}{{\color{red}\lantxh 高{\\ \bahao 可能增加腹泻、肠道炎症等风险}} }\end{center} \end{minipage} \\
% & \begin{minipage}{4.60cm}\begin{center}{{\color{red}\lantxh 高\\ \bahao 可能增加腹泻、肠道炎症等风险} }\end{center} \end{minipage} \\
\hline
\parbox[c]{\hsize}{\vskip7pt {\lantxh 嗜血杆菌属\\共生菌，某些菌种在一定条件下可引起化脓性感染、败血症等} \vskip7pt} & \parbox[c]{\hsize}{\vskip7pt\centerline{\raisebox{-1.5ex}{\includegraphics[width=5cm,height=0.55cm]{bar.pdf}}}\vskip7pt}  &
\hspace*{-1.51cm}\raisebox{-0.45ex}{\includegraphics{cry.pdf}}
 & \begin{minipage}{4.60cm}\begin{center}{{\color{red}\lantxh 高{\\ \bahao 可能增加化脓性感染、败血症的风险}} }\end{center} \end{minipage} \\
% & \begin{minipage}{4.60cm}\begin{center}{{\color{red}\lantxh 高\\ \bahao 可能增加化脓性感染、败血症的风险} }\end{center} \end{minipage} \\
\hline
\parbox[c]{\hsize}{\vskip7pt {\lantxh 柯林斯氏菌属\\发酵多种糖类产生氢气、乙醇等} \vskip7pt} & \parbox[c]{\hsize}{\vskip7pt\centerline{\raisebox{-1.5ex}{\includegraphics[width=5cm,height=0.55cm]{bar.pdf}}}\vskip7pt}  &
\hspace*{-4.83cm}\raisebox{-0.45ex}{\includegraphics{smile.pdf}}
 & \begin{minipage}{4.60cm}\begin{center}{{\lantxh 高{\\ \bahao 有利于发酵糖类}} }\end{center} \end{minipage} \\
% & \begin{minipage}{4.60cm}\begin{center}{{\color{green}\lantxh 高\\ \bahao 有利于发酵糖类} }\end{center} \end{minipage} \\
\hline
\parbox[c]{\hsize}{\vskip7pt {\lantxh 假单胞菌属\\分布广泛，部分菌种为条件致病菌，引起医院内获得性感染} \vskip7pt} & \parbox[c]{\hsize}{\vskip7pt\centerline{\raisebox{-1.5ex}{\includegraphics[width=5cm,height=0.55cm]{bar.pdf}}}\vskip7pt}  &
\hspace*{-1.51cm}\raisebox{-0.45ex}{\includegraphics{cry.pdf}}
 & \begin{minipage}{4.60cm}\begin{center}{{\color{red}\lantxh 高{\\ \bahao 可能增加感染风险}} }\end{center} \end{minipage} \\
% & \begin{minipage}{4.60cm}\begin{center}{{\color{red}\lantxh 高\\ \bahao 可能增加感染风险} }\end{center} \end{minipage} \\
\hline
\parbox[c]{\hsize}{\vskip7pt {\lantxh 韦荣氏球菌属\\分解乳酸为乙酸和丙酸，调节肠道pH，极少数菌种可能引起炎症} \vskip7pt} & \parbox[c]{\hsize}{\vskip7pt\centerline{\raisebox{-1.5ex}{\includegraphics[width=5cm,height=0.55cm]{bar.pdf}}}\vskip7pt}  &
\hspace*{-1.51cm}\raisebox{-0.45ex}{\includegraphics{cry.pdf}}
 & \begin{minipage}{4.60cm}\begin{center}{{\color{red}\lantxh 高{\\ \bahao 可能增加肠道炎症的风险}} }\end{center} \end{minipage} \\
% & \begin{minipage}{4.60cm}\begin{center}{{\color{red}\lantxh 高\\ \bahao 可能增加肠道炎症的风险} }\end{center} \end{minipage} \\
\hline
\parbox[c]{\hsize}{\vskip7pt {\lantxh 乳杆菌属\\肠道益生菌，能够生成乳酸，抑制有害菌及炎症，调节肠道环境} \vskip7pt} & \parbox[c]{\hsize}{\vskip7pt\centerline{\raisebox{-1.5ex}{\includegraphics[width=5cm,height=0.55cm]{bar.pdf}}}\vskip7pt}  &
\hspace*{-4.83cm}\raisebox{-0.45ex}{\includegraphics{smile.pdf}}
 & \begin{minipage}{4.60cm}\begin{center}{{\lantxh 高{\\ \bahao 有利于抑制肠道炎症及调节肠道环境}} }\end{center} \end{minipage} \\
% & \begin{minipage}{4.60cm}\begin{center}{{\color{green}\lantxh 高\\ \bahao 有利于抑制肠道炎症及调节肠道环境} }\end{center} \end{minipage} \\
\hline
\parbox[c]{\hsize}{\vskip7pt {\lantxh 肠球菌属\\多数为共生菌，产生乳酸。少数菌种可能导致多种感染症状} \vskip7pt} & \parbox[c]{\hsize}{\vskip7pt\centerline{\raisebox{-1.5ex}{\includegraphics[width=5cm,height=0.55cm]{bar.pdf}}}\vskip7pt}  &
\hspace*{-1.51cm}\raisebox{-0.45ex}{\includegraphics{cry.pdf}}
 & \begin{minipage}{4.60cm}\begin{center}{{\color{red}\lantxh 高{\\ \bahao 可能增加感染的风险}} }\end{center} \end{minipage} \\
% & \begin{minipage}{4.60cm}\begin{center}{{\color{red}\lantxh 高\\ \bahao 可能增加感染的风险} }\end{center} \end{minipage} \\
\hline
\parbox[c]{\hsize}{\vskip7pt {\lantxh 嗜胆菌属\\共生菌，可能与长期高脂高蛋白饮食有关} \vskip7pt} & \parbox[c]{\hsize}{\vskip7pt\centerline{\raisebox{-1.5ex}{\includegraphics[width=5cm,height=0.55cm]{bar.pdf}}}\vskip7pt}  &
\hspace*{-1.51cm}\raisebox{-0.45ex}{\includegraphics{cry.pdf}}
 & \begin{minipage}{4.60cm}\begin{center}{{\color{red}\lantxh 高{\\ \bahao 可能与长期高脂高蛋白饮食有关}} }\end{center} \end{minipage} \\
% & \begin{minipage}{4.60cm}\begin{center}{{\color{red}\lantxh 高\\ \bahao 可能与长期高脂高蛋白饮食有关} }\end{center} \end{minipage} \\
\hline
\parbox[c]{\hsize}{\vskip7pt {\lantxh 沙雷氏菌属\\多数为致病菌，导致尿路感染、腹膜炎等多种医院内感染} \vskip7pt} & \parbox[c]{\hsize}{\vskip7pt\centerline{\raisebox{-1.5ex}{\includegraphics[width=5cm,height=0.55cm]{bar.pdf}}}\vskip7pt}  &
\hspace*{-1.51cm}\raisebox{-0.45ex}{\includegraphics{cry.pdf}}
 & \begin{minipage}{4.60cm}\begin{center}{{\color{red}\lantxh 高{\\ \bahao 可能增加尿路感染、腹膜炎等的风险}} }\end{center} \end{minipage} \\
% & \begin{minipage}{4.60cm}\begin{center}{{\color{red}\lantxh 高\\ \bahao 可能增加尿路感染、腹膜炎等的风险} }\end{center} \end{minipage} \\
\hline
\parbox[c]{\hsize}{\vskip7pt {\lantxh 放线菌属\\多数为共生菌，少数在免疫力弱时可引起内源性感染、龋齿等} \vskip7pt} & \parbox[c]{\hsize}{\vskip7pt\centerline{\raisebox{-1.5ex}{\includegraphics[width=5cm,height=0.55cm]{bar.pdf}}}\vskip7pt}  &
\hspace*{-1.51cm}\raisebox{-0.45ex}{\includegraphics{cry.pdf}}
 & \begin{minipage}{4.60cm}\begin{center}{{\color{red}\lantxh 高{\\ \bahao 可能增加内源性感染、龋齿的风险}} }\end{center} \end{minipage} \\
% & \begin{minipage}{4.60cm}\begin{center}{{\color{red}\lantxh 高\\ \bahao 可能增加内源性感染、龋齿的风险} }\end{center} \end{minipage} \\
\hline
\parbox[c]{\hsize}{\vskip7pt {\lantxh 毛杆菌属\\肠道共生菌，发酵葡萄糖产生乳酸及少量乙酸和丁酸} \vskip7pt} & \parbox[c]{\hsize}{\vskip7pt\centerline{\raisebox{-1.5ex}{\includegraphics[width=5cm,height=0.55cm]{bar.pdf}}}\vskip7pt}  &
\hspace*{-4.83cm}\raisebox{-0.45ex}{\includegraphics{smile.pdf}}
 & \begin{minipage}{4.60cm}\begin{center}{{\lantxh 高{\\ \bahao 有利于互利共生及产生有益物质}} }\end{center} \end{minipage} \\
% & \begin{minipage}{4.60cm}\begin{center}{{\color{green}\lantxh 高\\ \bahao 有利于互利共生及产生有益物质} }\end{center} \end{minipage} \\
\hline
\parbox[c]{\hsize}{\vskip7pt {\lantxh 芽孢杆菌属\\多数为共生菌，降解复杂碳水化合物。少数可引起感染、食物中毒等} \vskip7pt} & \parbox[c]{\hsize}{\vskip7pt\centerline{\raisebox{-1.5ex}{\includegraphics[width=5cm,height=0.55cm]{bar.pdf}}}\vskip7pt}  &
\hspace*{-1.51cm}\raisebox{-0.45ex}{\includegraphics{cry.pdf}}
 & \begin{minipage}{4.60cm}\begin{center}{{\color{red}\lantxh 高{\\ \bahao 可能增加感染、食物中毒等的风险}} }\end{center} \end{minipage} \\
% & \begin{minipage}{4.60cm}\begin{center}{{\color{red}\lantxh 高\\ \bahao 可能增加感染、食物中毒等的风险} }\end{center} \end{minipage} \\
\hline
\parbox[c]{\hsize}{\vskip7pt {\lantxh 阿德勒氏菌属\\帮助分解食物中的异黄酮产生雌马酚，有利于肠道及人体健康} \vskip7pt} & \parbox[c]{\hsize}{\vskip7pt\centerline{\raisebox{-1.5ex}{\includegraphics[width=5cm,height=0.55cm]{bar.pdf}}}\vskip7pt}  &
\hspace*{-4.83cm}\raisebox{-0.45ex}{\includegraphics{smile.pdf}}
 & \begin{minipage}{4.60cm}\begin{center}{{\lantxh 高{\\ \bahao 有利于肠道及人体健康}} }\end{center} \end{minipage} \\
% & \begin{minipage}{4.60cm}\begin{center}{{\color{green}\lantxh 高\\ \bahao 有利于肠道及人体健康} }\end{center} \end{minipage} \\
\hline
\parbox[c]{\hsize}{\vskip7pt {\lantxh 克雷伯氏菌属\\多为致病菌，可能导致肺炎、尿路感染、软组织感染、菌血症等} \vskip7pt} & \parbox[c]{\hsize}{\vskip7pt\centerline{\raisebox{-1.5ex}{\includegraphics[width=5cm,height=0.55cm]{bar.pdf}}}\vskip7pt}  &
\hspace*{-1.51cm}\raisebox{-0.45ex}{\includegraphics{cry.pdf}}
 & \begin{minipage}{4.60cm}\begin{center}{{\color{red}\lantxh 高{\\ \bahao 可能增加肺炎、尿路感染、软组织感染、菌血症等风险}} }\end{center} \end{minipage} \\
% & \begin{minipage}{4.60cm}\begin{center}{{\color{red}\lantxh 高\\ \bahao 可能增加肺炎、尿路感染、软组织感染、菌血症等风险} }\end{center} \end{minipage} \\
\hline
\parbox[c]{\hsize}{\vskip7pt {\lantxh 脱硫弧菌属\\产生硫化氢，刺激肠道产生炎症反应，不利于肠道健康} \vskip7pt} & \parbox[c]{\hsize}{\vskip7pt\centerline{\raisebox{-1.5ex}{\includegraphics[width=5cm,height=0.55cm]{bar.pdf}}}\vskip7pt}  &
\hspace*{-1.51cm}\raisebox{-0.45ex}{\includegraphics{cry.pdf}}
 & \begin{minipage}{4.60cm}\begin{center}{{\color{red}\lantxh 高{\\ \bahao 可能增加肠道炎症风险}} }\end{center} \end{minipage} \\
% & \begin{minipage}{4.60cm}\begin{center}{{\color{red}\lantxh 高\\ \bahao 可能增加肠道炎症风险} }\end{center} \end{minipage} \\
\hline
\parbox[c]{\hsize}{\vskip7pt {\lantxh 爱格士氏菌属\\多为致病菌，与溃疡性结肠炎、肛门脓肿、菌血症等疾病有关} \vskip7pt} & \parbox[c]{\hsize}{\vskip7pt\centerline{\raisebox{-1.5ex}{\includegraphics[width=5cm,height=0.55cm]{bar.pdf}}}\vskip7pt}  &
\hspace*{-1.51cm}\raisebox{-0.45ex}{\includegraphics{cry.pdf}}
 & \begin{minipage}{4.60cm}\begin{center}{{\color{red}\lantxh 高{\\ \bahao 可能增加溃疡性结肠炎、肛门脓肿、菌血症等的风险}} }\end{center} \end{minipage} \\
% & \begin{minipage}{4.60cm}\begin{center}{{\color{red}\lantxh 高\\ \bahao 可能增加溃疡性结肠炎、肛门脓肿、菌血症等的风险} }\end{center} \end{minipage} \\
\hline
\parbox[c]{\hsize}{\vskip7pt {\lantxh 乳球菌属\\发酵葡萄糖产生乳酸，调节肠道pH，抑制有害菌的生长} \vskip7pt} & \parbox[c]{\hsize}{\vskip7pt\centerline{\raisebox{-1.5ex}{\includegraphics[width=5cm,height=0.55cm]{bar.pdf}}}\vskip7pt}  &
\hspace*{-4.83cm}\raisebox{-0.45ex}{\includegraphics{smile.pdf}}
 & \begin{minipage}{4.60cm}\begin{center}{{\lantxh 高{\\ \bahao 有利于调节肠道pH及抑制有害菌生长}} }\end{center} \end{minipage} \\
% & \begin{minipage}{4.60cm}\begin{center}{{\color{green}\lantxh 高\\ \bahao 有利于调节肠道pH及抑制有害菌生长} }\end{center} \end{minipage} \\
\hline
\parbox[c]{\hsize}{\vskip7pt {\lantxh 真杆菌属\\产生短链脂肪酸等有益物质，少数可能与炎症、菌血症等相关} \vskip7pt} & \parbox[c]{\hsize}{\vskip7pt\centerline{\raisebox{-1.5ex}{\includegraphics[width=5cm,height=0.55cm]{bar.pdf}}}\vskip7pt}  &
\hspace*{-1.51cm}\raisebox{-0.45ex}{\includegraphics{cry.pdf}}
 & \begin{minipage}{4.60cm}\begin{center}{{\color{red}\lantxh 高{\\ \bahao 可能增加肠道炎症、菌血症等的风险}} }\end{center} \end{minipage} \\
% & \begin{minipage}{4.60cm}\begin{center}{{\color{red}\lantxh 高\\ \bahao 可能增加肠道炎症、菌血症等的风险} }\end{center} \end{minipage} \\
\hline
\parbox[c]{\hsize}{\vskip7pt {\lantxh 埃希氏菌属\\多数为共生菌，产生维生素K，少数与腹泻、尿道感染、败血症等相关} \vskip7pt} & \parbox[c]{\hsize}{\vskip7pt\centerline{\raisebox{-1.5ex}{\includegraphics[width=5cm,height=0.55cm]{bar.pdf}}}\vskip7pt}  &
\hspace*{-1.51cm}\raisebox{-0.45ex}{\includegraphics{cry.pdf}}
 & \begin{minipage}{4.60cm}\begin{center}{{\color{red}\lantxh 高{\\ \bahao 可能增加腹泻、尿道感染、败血症等的风险}} }\end{center} \end{minipage} \\
% & \begin{minipage}{4.60cm}\begin{center}{{\color{red}\lantxh 高\\ \bahao 可能增加腹泻、尿道感染、败血症等的风险} }\end{center} \end{minipage} \\
\hline
\parbox[c]{\hsize}{\vskip7pt {\lantxh 别样杆菌属\\共生菌，降解多糖产生琥珀酸和乙酸，可能与长期高脂饮食相关} \vskip7pt} & \parbox[c]{\hsize}{\vskip7pt\centerline{\raisebox{-1.5ex}{\includegraphics[width=5cm,height=0.55cm]{bar.pdf}}}\vskip7pt}  &
\hspace*{-1.51cm}\raisebox{-0.45ex}{\includegraphics{cry.pdf}}
 & \begin{minipage}{4.60cm}\begin{center}{{\color{red}\lantxh 高{\\ \bahao 可能与长期高脂饮食有关}} }\end{center} \end{minipage} \\
% & \begin{minipage}{4.60cm}\begin{center}{{\color{red}\lantxh 高\\ \bahao 可能与长期高脂饮食有关} }\end{center} \end{minipage} \\
\hline
\parbox[c]{\hsize}{\vskip7pt {\lantxh 丁酸弧菌属\\降解多种碳水化合物，产生丁酸，参与蛋白质和脂类代谢} \vskip7pt} & \parbox[c]{\hsize}{\vskip7pt\centerline{\raisebox{-1.5ex}{\includegraphics[width=5cm,height=0.55cm]{bar.pdf}}}\vskip7pt}  &
\hspace*{-4.83cm}\raisebox{-0.45ex}{\includegraphics{smile.pdf}}
 & \begin{minipage}{4.60cm}\begin{center}{{\lantxh 高{\\ \bahao 有利于产生有益物质及物质代谢}} }\end{center} \end{minipage} \\
% & \begin{minipage}{4.60cm}\begin{center}{{\color{green}\lantxh 高\\ \bahao 有利于产生有益物质及物质代谢} }\end{center} \end{minipage} \\
\hline
\parbox[c]{\hsize}{\vskip7pt {\lantxh 弯曲杆菌属\\多数菌种为致病菌，可引起弯曲菌病，表现为严重腹泻或痢疾综合征} \vskip7pt} & \parbox[c]{\hsize}{\vskip7pt\centerline{\raisebox{-1.5ex}{\includegraphics[width=5cm,height=0.55cm]{bar.pdf}}}\vskip7pt}  &
\hspace*{-1.51cm}\raisebox{-0.45ex}{\includegraphics{cry.pdf}}
 & \begin{minipage}{4.60cm}\begin{center}{{\color{red}\lantxh 高{\\ \bahao 可能增加腹泻或痢疾风险}} }\end{center} \end{minipage} \\
% & \begin{minipage}{4.60cm}\begin{center}{{\color{red}\lantxh 高\\ \bahao 可能增加腹泻或痢疾风险} }\end{center} \end{minipage} \\
\hline
\multicolumn{3}{l}{}\\
\multicolumn{3}{l}{\includegraphics[scale=1]{greencircle.pdf} {\fontsize{9.3pt}{11pt}\selectfont\bf 绿色表示健康}   \includegraphics[scale=1]{orangecircle.pdf}  {\fontsize{9.3pt}{11pt}\selectfont\bf 橙色表示需要关注}  \includegraphics[scale=1]{redcircle.pdf} {\fontsize{9.3pt}{11pt}\selectfont\bf 红色表示有风险}}
\end{longtable}

\vspace*{1mm}
\begin{LRaside}[.8]{结果分析}
\noindent
\includegraphics[width=\linewidth]{result_fenbu.pdf}
\asidebreak %
综合您的肠道菌群检测结果，
您肠道内的拟杆菌属、多尔氏菌属、链球菌属、普雷沃氏菌属、萨特氏菌属等指标异常，可能增加肠易激综合征、化脓性炎症、肺炎等的风险，不利于肠道菌群平衡。您体内的布劳特氏菌属、栖粪杆菌属、粪球菌属、瘤胃球菌属、颤螺菌属等指标优于参考人群，有利于产生有益物质、肠道健康、抑制肠道炎症等。肠道菌群含量维持动态平衡，少数菌种数量的改变并非致病的直接原因，整体肠道菌群失衡，才可能引起疾病的发生。
\end{LRaside}

\end{document}

